\subsection{权群的情形；极大项}

保留前一号的记号，并令 R 为实向量空间 V 中的约化根系。在本节余下部分，取 P 为 R 的权群（\(\S\) 1，第 9 号）。群 W = W(R) 作用于 P，因而也作用于代数 A{[}P{]}；有 \(w(e^{p}) = e^{w(p)}\)（其中 \(w \in W\) 且 \(p \in P\)）。

令 C 为 R 的一个室（\(\S\) 1，第 5 号），令 \(B = (\alpha_{i})_{1 \leqslant i \leqslant l}\) 为相应的 R 的一组基。赋予 V（从而也赋予 P）由 C 定义的序结构。若 \(p, p' \in P\)，则 \(p \geqslant p'\) 当且仅当 \(p - p'\) 是具有正系数的 \(\alpha_{i}\) 的线性组合。

定义 1. 设 \(x = \sum_{p \in \mathrm{P}} x_p e^p\) 是 A{[}P{]} 的一个元素。令 S 为满足 \(p \in \mathrm{P}\) 且 \(x_p \neq 0\) 的指标集合，称 S 为 \(x\) 的支集；令 X 为 S 的极大元集合，称 X 为 \(x\) 的最大支集。于是，称 \(x_p e^p\)（其中 \(p \in X\)）为 \(x\) 的极大项。

引理 2. 设 \(x \in \mathrm{A}[\mathrm{P}]\)，且 \((x_{p}e^{p})_{p \in \mathrm{X}}\) 是 \(x\) 的极大项族。设 \(q \in \mathrm{P}\)，且 \(y \in \mathrm{A}[\mathrm{P}]\) 满足 \(e^{q}\) 是 \(y\) 的唯一极大项。则 \(xy\) 的极大项族为 \((x_{p}e^{p + q})_{p \in \mathrm{X}}\)。

置 \(x = \sum_{p} x_{p} e^{p}, y = \sum_{r} y_{r} e^{r}\) 且 \(xy = \sum_{t} z_{t} e^{t}\)。则 \(r \leqslant q\)（对所有 \(r \in \mathbb{P}\) 且满足 \(y_{r} \neq 0\) 的情形），并且 \(z_{t} = \sum_{p + r = t} x_{p} y_{r}\)。

若 \(t = p + q = p' + r\)，其中 \(p \in \mathrm{X}\) 且 \(x_{p'}y_r \neq 0\)，则 \(r \leqslant q\)，故 \(p' \geqslant p\)，从而 \(p' = p\)。于是 \(z_{p + q} = x_p y_q = x_p \neq 0\)。这表明 \(\mathrm{X} + q\) 包含于乘积 \(xy\) 的支集。

另一方面，若 \(t = p' + r\) 且 \(x_{p'} y_r \neq 0\)，则存在 \(p \in X\) 使得 \(p' \leqslant p\)，于是 \(t \leqslant p + q\)。因此 xy 的最大支集包含于 \(X + q\)。由于 \(X + q\) 中任意两个元素都不可比较，故 \(X + q\) 恰好是 xy 的最大支集，并且如上所见，有 \(z_{p+q} = x_p\)（对 \(p \in X\)），引理得证。

备注。由于 \(x \neq 0\) 意味着 x 的最大支集非空，引理 2 表明，\(x \neq 0\) 蕴含 \(xy \neq 0\)，只要 y 有形如 \(e^{q}\) 的唯一极大项。

